\section{The Abrigo Log: Record of Closed and Active Iterations}\label{sec:6}

This section records the five iterations the project has run to date.
Each iteration is presented in the same nine-field schema --- target
population; \((Y, X)\) pair with operational definition; sign
expectation and literature anchor; primary specification; sample;
verdict; headline numbers; key caveats from independent review; and
pre-registration sha256 chain. The section closes with a load-bearing
caveat on cross-iteration multiple-comparison structure that we develop
further as a supervisor question in \(\S\ref{sec:7}\) (\(\S 7.5\)).

\subsection{Pair D --- BPO offshoring through COP/USD lag}\label{sec:6.1}

\paragraph{Target population.} Colombian young workers aged 14--28
(statutory youth band per Ley~1622 de 2013) employed in the broad
services aggregate.

\paragraph{(Y, X) pair.}
\(Y_t\) is the share of young workers (14--28) whose primary occupation
falls in CIIU Rev.~4~A.C. Sections G--T, computed monthly on the DANE
GEIH micro-data with \texttt{FEX\_C\_2018} expansion factors and the
DANE \emph{empalme} factor applied to the 2015--2020 segment for
Marco-2005-to-Marco-2018 reconciliation; the dependent variable is the
logit of the raw share. \(X_{t-k}\) is \(\log(\mathrm{COP\_USD}_{t-k})\)
on the Banrep TRM end-of-month spot, for \(k \in \{6, 9, 12\}\) months.

\paragraph{Sign expectation and literature anchor.} Positive. The
transmission chain is Baumol cost disease
\citep{baumolbowen1966, baumol2006} composed with US-Colombia BPO
wage arbitrage \citep{mendietamunoz2017} and the standard
six-to-twelve-month offshoring contracting cycle documented in the
Philippines comparator \citep{beerepoothendriks2013, errighibodwell2017};
displaced labor reallocates into services on the Rodrik-McMillan
premature-deindustrialization channel \citep{rodrik2016, mcmillanrodrik2011}.

\paragraph{Primary specification.} Equation~(\ref{eq:primary-monthly})
with composite \(\beta_{\mathrm{composite}} = \beta_6 + \beta_9 +
\beta_{12}\); one-sided test at \(\alpha = 0.05\); OLS-homoskedastic
inference primary, HAC(\(L = 12\)) inference as robustness arm R4.

\paragraph{Sample.} \(N = 134\) monthly observations, window
2015-01-31 through 2026-02-28. The window's lower bound is the result
of two CORRECTIONS-block revisions documented in the spec: an initial
2008-01 lower bound was shortened to 2010-01 after a schema-stability
HALT on the DANE empalme nota-técnica coverage, and then to 2015-01
after a second HALT under the realization that the empalme files for
the 2010--2014 segment ship CIIU Rev.~3 codes in the
\texttt{RAMA4D} column despite the column header naming. The
2015-01 lower bound is the DANE-canonical Rev.~4 value-content-verified
window. The \(N = 134\) realized sample is well above the
\(N_{\min} = 75\) floor.

\paragraph{Verdict.} PASS, 2026-04-28.

\paragraph{Headline numbers.} \(\hat{\beta}_{\mathrm{composite}} =
+0.13670985\); HAC standard error \(0.02465\); \(t = +5.5456\);
one-sided \(p \approx 1.46 \times 10^{-8}\). Robustness rows R1--R4
produced zero sign flips relative to the primary; the four-row classifier
returned AGREE.

\paragraph{Key caveats from independent review.} Five flags from the
Reality Checker review of the closed iteration are inherited as
load-bearing constraints on any downstream framing:
\begin{enumerate}
  \item The estimate hedges the \emph{correlation} between lagged COP/USD
    and the broad-services young-worker share. It does not identify the
    BPO transmission channel as the operative mechanism: any narrative
    causal interpretation of the BPO chain is conjectural and is not
    supported by the Stage-1 design.
  \item The lag effect is concentrated at the six-month horizon
    (\(\hat{\beta}_6 \approx +0.109\), approximately 80\% of the
    composite) rather than uniformly spread across the 6--12-month
    contracting window. The narrative that the project carries forward
    is ``concentrated at the six-month horizon, within the
    six-to-twelve-month window,'' not ``spread uniformly across
    the window.''
  \item The window-narrowing path through the two CORRECTIONS blocks
    over-represents three macroeconomically distinctive regimes
    (post-2014 oil-price collapse, the COVID shock, and the 2022--2024
    Federal-Reserve tightening cycle); the regime-mix concern is
    a residual identification concern beyond the four pre-registered
    robustness arms.
  \item The verdict is sensitive to a single off-spec design choice. An
    orchestrator-brief variant that included a \texttt{marco2018\_dummy}
    in the primary --- not in the spec text, where the dummy is the R1
    arm only --- produces \(\hat{\beta}_{\mathrm{composite}} = +0.0815\)
    at one-sided \(p \approx 0.080\), which would map to ESCALATE
    Clause~A under the verdict tree of \(\S\ref{sec:5.6}\). The spec
    text ran as authoritative because the spec sha256 was pinned
    9.5 hours before Phase-2 dispatch; the off-spec variant is preserved
    in the artifact \texttt{primary\_ols.json} under the field
    \texttt{off\_spec\_sensitivity\_orchestrator\_brief} for audit. The
    document carries this verdict-sensitivity record forward as part of
    the framework's decision log.
  \item \emph{Compositional reframe surfaced by the dev-AI Stage-1
    closed iteration (2026-05-06).} The Pair~D PASS was estimated on
    the broad CIIU Sections G--T cut and was, at spec-pin time,
    interpreted under a ``broad-services hedging the BPO transmission
    chain'' reading consistent with either a Section-J ICT-narrow
    mechanism or a Section-M-style consultancy / scientific R\&D /
    administrative-services mechanism aggregated up. The dev-AI
    Stage-1 iteration's Section-J primary closed
    \(\hat{\beta}_{\mathrm{composite}} = -0.14613\)
    (HAC(\(L=12\)) SE \(0.0847\); \(t = -1.726\); one-sided
    \(p \approx 0.958\) on the pre-registered \(H_1: \beta > 0\)) ---
    a sign-flipped FAIL --- and its R2 sensitivity arm on Section M
    closed \(\hat{\beta}_{\mathrm{composite}} = +0.45482801\)
    (\(t = +4.73\); \(p = 1.13 \times 10^{-6}\)) at the same monthly
    cadence on the same lag panel. Read jointly with Pair~D, this
    pattern strongly suggests Pair~D's PASS was carried by the
    Section-M-style sub-aggregate (consultants, legal, accounting,
    scientific R\&D, administrative services) rather than by Section-J
    ICT-narrow tech offshoring. The spec~\(\S 9.16\) compositional-
    accounting concern is therefore \emph{empirically resolved as
    flagged-not-resolved}: a formal (Sections G--T \(\setminus\) J)
    decomposition pre-registered as an R5 robustness row remains
    deferred for a future iteration, and any Stage-2 dispatch off the
    Pair~D PASS must inherit this compositional caveat. Pair~D's
    headline numbers stand verbatim under the original spec; the
    interpretation of the operative sub-aggregate is what has
    sharpened.
\end{enumerate}

\paragraph{Pre-registration sha256 chain.}
v1.1 \(\to\) v1.2 \(\to\) v1.2.1 \(\to\) v1.3 \(\to\) v1.3.1
\(\to\) governing hash
\texttt{964c62cca0...ef659}. The full hash chain and revision history
are reproduced in Appendix~\ref{app:C}.

\subsection{dev-AI Stage-1 --- Section J ICT through COP/USD lag}\label{sec:6.2}

\paragraph{Target population.} Colombian young workers (14--28) employed
in CIIU Rev.~4 A.C. Section J (Information and Communication, Divisions
58--63: publishing, motion picture and video, programming and
broadcasting, telecommunications, computer programming and consultancy,
information service activities). The population is a strict subset of
the Pair~D Sections G--T cut.

\paragraph{(Y, X) pair.} \(Y_t\) is the Section-J young-worker share
computed on DANE GEIH micro-data with the same monthly cadence,
\texttt{FEX\_C\_2018} expansion factors, and \texttt{RAMA4D\_R4}
sector coding as Pair~D, restricted to Section J via the published
DANE Section table; the dependent variable is the logit. \(X_{t-k}\)
is the same Banrep TRM end-of-month-spot panel as in Pair~D, with the
same lag set \(k \in \{6, 9, 12\}\); the panel back-extends \(X\) to
2014-01 to preserve \(Y\) starting at 2015-01.

\paragraph{Sign expectation and literature anchor.} Positive. The
mechanism was the same Baumol-arbitrage-offshoring chain as Pair~D
narrowed from BPO-flavored services-broad to ICT-narrow
tech-services offshoring; the same canonical references applied
\citep{mendietamunoz2017, beerepoothendriks2013, errighibodwell2017,
rodrik2016}. As recorded below, the sign expectation was
\emph{rejected at conventional significance} by the realized estimate.

\paragraph{Primary specification.} Equation~(\ref{eq:primary-monthly})
restated with the Section-J \(Y\); one-sided test at \(\alpha = 0.05\);
OLS-homoskedastic inference primary, HAC(\(L = 12\)) inference as
robustness arm R4.

\paragraph{Sample.} \(N = 134\) monthly observations, window
2015-01-31 through 2026-03-31, well above the \(N_{\min} = 75\) floor.

\paragraph{Verdict.} CLOSED FAIL, 2026-05-06.

\paragraph{Headline numbers.}
\(\hat{\beta}_{\mathrm{composite}} = -0.14613\) (HAC(\(L=12\))
standard error \(0.0847\)); \(t = -1.726\); one-sided
\(p \approx 0.958\) against \(H_1: \beta > 0\). The estimated sign is
the opposite of the pre-registered expectation: the FAIL is a
\emph{sign-flipped} FAIL, not an underpowered failure-to-reject.
The R-consistency classifier returned MIXED (three of four R-rows agree
with the primary; only the R2 Section-M sensitivity sign-flips). The
v1.0.2 pre-registered hardening (R1 \emph{and} R3 must individually
agree on sign for the verdict to clear PASS independently of the
four-row aggregate) is consistent with the FAIL verdict at the
negative sign. The pre-authorized escalation suite (\(\S\ref{sec:5.6}\))
returned: D-i and D-ii FAIL on the positive-sign criterion; D-iii
literally returned \(\hat{\beta} = +0.113\), \(p = 0.012\) but is
\emph{inadmissible} per the strict reading of \(\S\ref{sec:5.5}\)
because the \(\S 3.3\)-trigger condition that would license escalation
never fired. The user closed FAIL on 2026-05-06 per the strict
\(\S\ref{sec:5.5}\) reading, with the D-iii literal positive preserved
on the record as a Phase-3 artifact for a possible future Section-M
iteration. (A spec-internal contradiction between
\(\S 5.5\)~line~252 and \(\S 9.6\) on the question of whether the
escalation suite must run conditional on a primary-trigger or
unconditionally is acknowledged on the project's deferred-spec-
reconciliation list as \emph{spec~v1.0.3 reconciliation flag},
to be resolved unambiguously before another \(\S 5.5\) invocation.)

\paragraph{Key caveats from independent review.}
\begin{enumerate}
  \item \emph{Cell-count and share-magnitude gap (CORRECTIONS-\(\kappa\)
    block).} Realized monthly cell counts for the Section-J young-worker
    stratum lay in \([94, 267]\) with median \(145\); the ex-ante
    feasibility-memo baseline was \([700, 1200]\). Realized monthly raw
    shares lay in \([0.014, 0.031]\) against an ex-ante range of
    \([0.04, 0.10]\). The logit derivative
    \(\mathrm{d}\,\mathrm{logit}(Y)/\mathrm{d}Y = 1/[Y(1-Y)]\) at the
    realized share range maps to \([33, 73]\), against a v1.0.1
    across-support amplification ratio of \(2.34\); the spec was
    \(\kappa\)-tightened on this basis at v1.0.2 to require R1 and R3
    individual sign agreement before any PASS verdict. The realized
    estimate cleared the \(\kappa\)-tightened reading at the
    \emph{negative} sign --- the substrate is not too noisy under the
    DISAGREE classifier (only one R-row sign-flips), so the verdict is
    routable rather than SUBSTRATE\_TOO\_NOISY.
  \item \emph{R2 Section-M sensitivity --- compositional finding.}
    The R2 sensitivity arm narrowed \(Y\) from Section J to Section M
    (CIIU Rev.~4 Sections 69--75: professional, scientific, technical,
    and administrative services) at the same lag panel and primary
    functional form, and produced
    \(\hat{\beta}_{\mathrm{composite}} = +0.45482801\) at
    \(t = +4.73\), one-sided \(p = 1.13 \times 10^{-6}\). This is
    a strong positive signal at a sister sub-aggregate to Section J
    on the same FX-lag panel. The framework reads this as a
    \emph{candidate-next-iteration finding}, not as a Pair~D rescue or
    a dev-AI Stage-1 rescue: the Section-M coefficient was not
    pre-registered as the primary at spec-pin time, and per the
    project's anti-fishing discipline it cannot graduate to a
    PASS-equivalent claim from a robustness-arm posture. What it
    \emph{does} resolve, jointly with the Pair~D PASS and the dev-AI
    Stage-1 FAIL, is the spec~\(\S 9.16\) compositional-accounting
    ambiguity: Pair~D's broad-services PASS is \emph{empirically
    consistent with} a Section-M-style consultancy-and-professional-
    services transmission chain rather than a Section-J ICT-narrow
    chain. A future iteration with Section M pre-registered as the
    primary, on a fresh spec-pin and the canonical four-row R-arm
    discipline, is the routed-forward path.
  \item \emph{Empalme residual bias surfaced.} The Phase-1 ingest
    diagnostic (Notebook 02 Trio 1 \texttt{boundary\_anomaly}) returned
    TRUE: the logit-\(Y\) jump at the 2020-12 \(\to\) 2021-01 boundary
    was \(+0.375\), against an envelope of \(\pm 0.335\) (three times
    the historical month-to-month standard deviation). The R1
    regime-dummy arm produced \(\hat{\beta}_{\mathrm{R1}} = +0.188\)
    at \(t = +4.36\), reflecting an unexplained level shift in the
    post-2021 era. The DANE Marco-2018 empalme correction was applied
    to the primary as designed, but did \emph{not} fully neutralize the
    methodology break at the Section-J narrow cut; the regime structure
    is informational diagnostic and the project explicitly does not
    use the \(\hat{\beta}_{\mathrm{R1}}\) sign as a primary-spec
    rescue lever. The empalme-residual phenomenon is an inherited
    methodology constraint for any future DANE-GEIH-based iteration.
  \item \emph{Per-user instrument design implication.} The Stage-1
    rejection of the positive-sign expectation falsifies the
    transmission for which a per-user a\(_s\) instrument was the
    candidate Stage-2 deployment: the rejected sign means the dev-AI-
    paying population (Colombian young Section-J workers paying
    USD-denominated AI APIs) does not, at Stage-1, exhibit the
    FX-pass-through-to-employment pattern the convex hedge would
    settle. A SHORT-FX-style instrument would be the opposite of what
    that population needs given USD-denominated AI tooling costs;
    consequently a per-user a\(_s\) construction on this transmission
    is not designable. The Section-M finding does not change this
    conclusion at Section-J granularity, and any Section-M follow-on
    iteration must pre-pin its own population definition before the
    Stage-2 instrument question reopens.
\end{enumerate}

\paragraph{Pre-registration sha256 chain.} v1.0 \(\to\) v1.0.1
\(\to\) v1.0.2 \(\to\) governing hash
\texttt{7c72292516...751f5a}.

\subsection{FX-vol-CPI surprise --- Colombian CPI through TRM realised
volatility}\label{sec:6.3}

\paragraph{Target population.} Colombian corporates and remittance
recipients with FX-volatility exposure on the COP/USD pair around
Colombian CPI publication.

\paragraph{(Y, X) pair.} \(Y_t\) is the weekly realized volatility of
the Banrep TRM end-of-day series, computed on the daily log-return
panel and aggregated to the weekly cadence under a Friday anchor.
\(X_t\) is the signed AR(1) CPI surprise: the monthly DANE CPI
release minus a one-step-ahead AR(1) forecast on prior CPI prints,
mapped to the weekly cadence by carrying the surprise across the
release week and zeros otherwise.

\paragraph{Sign expectation and literature anchor.} Positive on the
classic CPI-announcement-volatility channel: a larger absolute
surprise should generate larger TRM realized volatility around the
release. The pre-registered primary tested the signed surprise rather
than the absolute surprise, on the prior that the Colombian
hyperinflation anchor produces an asymmetric response.

\paragraph{Primary specification.} Weekly OLS of TRM realized
volatility on the signed AR(1) CPI surprise plus a small set of
macro-surprise controls; one-sided test at \(\alpha = 0.05\);
OLS-homoskedastic inference primary, HAC(\(L = 4\)) inference under
the gate-aggregation rule. The pre-registered sensitivity universe
spans thirteen rows (A1--A12 plus seven cumulative cuts S1--S7)
hashed via a \texttt{sensitivity\_preregistration\_hash} on the
panel fingerprint artifact.

\paragraph{Sample.} \(N = 947\) weekly observations, window
2008-01-02 through 2026-02-23.

\paragraph{Verdict.} CLOSED FAIL on the pre-registered weekly primary,
2026-04-19.

\paragraph{Headline numbers.} \(\hat{\beta}_{\mathrm{CPI}} = -0.000685\)
(HAC(4) standard error \(0.001794\)); \(90\%\) HAC interval
\([-0.003636,\,+0.002266]\) (contains zero); the gate criterion
\(\hat{\beta} - 1.28 \cdot \widehat{\mathrm{SE}} < 0\) returns
\(-0.002981 < 0\), which is the FAIL trigger. Adjusted \(R^2 = 0.193\)
clears the pre-pinned floor; the macro-surprise regressors alone
contribute approximately zero. Under the pre-committed weekly OLS
specification, Colombian AR(1) CPI surprises do not produce a
statistically detectable increase in COP/USD realized volatility
across the sample window. We separately note that the pre-registered
A1 monthly-cadence sensitivity returned
\(\hat{\beta} = +0.0152\) with a \(90\%\) interval
\([+0.0057,\,+0.0246]\), and the A4 release-day-excluded sensitivity
returned \(\hat{\beta} = +0.0033\) with a \(90\%\) interval
\([+0.0005,\,+0.0062]\). The A4 \(90\%\) interval excludes zero
uncorrected: the rejection at the one-sided \(\alpha = 0.05\)
convention is clean before any multiple-testing correction, even
though the lower bound (\(0.0005\)) is small in absolute magnitude;
calling A4 ``edge-of-significance'' would be inaccurate. Neither A1
nor A4 is the iteration's main finding: the pre-registered primary is
FAIL and the pre-registered primary \emph{is} the result. The two
sensitivities are pre-registered robustness checks that returned
positive significance at the \(90\%\) one-sided convention; they live
in the iteration's pivot record as candidate inputs to a future
monthly-cadence-primary specification, not as a rescue of the
weekly primary.

\paragraph{Key caveats from independent review.} The pre-registered
T1 exogeneity test rejected at \(F = 15.12\),
\(p \approx 1.3 \times 10^{-9}\); the iteration is therefore a
\emph{predictive regression} in the strict sense of the
Granger-predictive-versus-structural-causal distinction, and the
estimated \(\hat{\beta}\) describes how past information predicts
current realized volatility, not how an exogenous surprise causes
future realized volatility. We do not interpret \(\hat{\beta}\) as an
impulse-response. The HAC-versus-bootstrap reconciliation across the
verdict tree returned AGREEMENT, ruling out a standard-error artifact.

\paragraph{Pre-registration sha256 chain.} The iteration ran on the
Phase-A.0 Rev~4 specification chain; the canonical artifacts are the
panel fingerprint
(\texttt{nb1\_panel\_fingerprint.json}) and the verdict file
(\texttt{gate\_verdict.json}), both governed by the
\texttt{sensitivity\_preregistration\_hash} pinned in the panel
fingerprint.

\subsection{Phase-A.0 remittance --- Colombian remittance corridor through
FX}\label{sec:6.4}

\paragraph{Target population.} Colombian household remittance recipients,
indexed against the Colombian remittance corridor with FX pass-through.

\paragraph{(Y, X) pair.} \(Y_t\) was the weekly TRM realized volatility
panel (inherited from \(\S\ref{sec:6.3}\)). \(X_t\) was a candidate
on-chain remittance-flow proxy constructed from cCOP and COPm aggregate
user activity on Celo, partitioned to attempt to isolate Colombian
household-remittance flow from third-party-DEX activity, treasury
roundtripping, and bot arbitrage.

\paragraph{Sign expectation and literature anchor.} Positive on the
remittance-volatility channel: large remittance-flow shocks should
generate FX-vol shocks at the receiving corridor. The literature anchor
was the parametric-insurance-counterparty-matching tradition.

\paragraph{Primary specification.} Pre-registered weekly OLS with
empirical kill-criteria k1 (intent: an X-driver fingerprint must be
present in remittance event-window data) and k2 (data-source: the X
proxy must measure Colombian household remittance, not a generic
on-chain volume aggregate).

\paragraph{Sample.} Window 2024-09 through 2026-04, weekly cadence
(panel construction did not reach final \(N\) before the EXIT verdict
fired).

\paragraph{Verdict.} CLOSED EXIT\_NON\_REMITTANCE, 2026-04-24.

\paragraph{Headline numbers.} The iteration exited at the kill-criteria
step rather than producing point estimates. Task~11.F Axis-1
event-research returned zero of thirty peak-day flows that fingerprinted
as remittance; approximately 87\% of the surveyed flow decomposed into
TRM arbitrage, treasury roundtripping, campaigns, and bot activity, with
the residual a thin retail tail. Kill criterion k1 fired (intent
non-remittance dominant); kill criterion k2 fired in part (the
unfiltered Dune aggregate at query 7366593 did not measure Colombian
household remittance).

\paragraph{Key caveats from independent review.} The exit was load-bearing
on the kill-criterion discipline. Without the pre-registered EXIT
criterion, the rev-4 plan's ``non-stop filter iteration'' policy would
have run to a filter that appeared to ``rescue'' a remittance signal
that was never captured upstream. The iteration carries forward the
lesson that no post-hoc filter can reconstruct what was never measured
upstream; any pivot must inherit the kill-criterion pattern.

A further caveat applies to subsequent on-chain attribution. A
2026-04-24 thesis proposed Carbon-DeFi basket-rebalancing volume as a
candidate on-chain demand signal for an inequality-differential
follow-on iteration; that attribution was retracted on 2026-04-27 after
verification of the Mento V3 deployment manifest confirmed zero
protocol-level integration between the Mento contracts and the Carbon
contracts. Carbon \texttt{tokenstraded} events on Celo are pure
third-party-DEX activity dominated by Bancor's own arbitrage routers
(\(383{,}303\) distinct traders on the Mento Broker V2 swap event versus
\(147\) on Carbon over comparable windows), not Mento Reserve user
demand. The framework's on-chain demand-signal candidate has been
re-pivoted to native Mento Broker V2 swap events on the Broker contract
\texttt{0x777A8255...4CaD}, and the supervisor-review document does
\emph{not} present Carbon-DeFi volume as the canonical Phase-A.0
finding. The Phase-A.0 EXIT verdict stands on the original
\((Y, X)\) pair as falsified-on-mechanism rather than failed-on-power.

\paragraph{Pre-registration sha256 chain.} Phase-A.0 spec chain
Rev~1 \(\to\) Rev~1.1 \(\to\) Rev~1.1.1 (all SUPERSEDED and RETIRED at
EXIT) \(\to\) Rev~4.1 (kill-criterion-bearing); the on-chain attribution
retraction is recorded in the framework memory at Rev-5.3.7.

\subsection{P1 Bittensor SN18 --- event-study apparatus}\label{sec:6.5}

\paragraph{Target population.} The proprietary-AI policy-shock-exposed
population (LATAM developers paying USD-denominated AI APIs is the
downstream candidate; the apparatus was authored before the
target-population narrowing).

\paragraph{(Y, X) pair.} \(Y_t\) was a Bittensor SN18 (Cortex.t) alpha
event-study indicator; \(X_t\) was a set of pre-pinned policy-event
dates (Bittensor halving 2025-12-15; dTAO mainnet activation 2025-02-13;
Cortex.t milestone events).

\paragraph{Sign expectation and literature anchor.} Two-sided. The
methodology baseline was the Maymin (2026, arXiv:2603.29751)
event-study framework with explicit divergences pinned in the spec.

\paragraph{Primary specification.} Event-study apparatus over a
nine-cell verdict cube (verdict-eligibility crossed with
robustness-consistency), Bonferroni-corrected at
\(\alpha_{\mathrm{primary}} = 0.0167\) over the conjunctive family
P1 + P2 + P3, with \(N_{\min}^{\mathrm{events}} = 8\) and an
asymmetric placebo gate.

\paragraph{Sample.} Apparatus completed; data ingest deferred.

\paragraph{Verdict.} PARKED, 2026-04-27.

\paragraph{Headline numbers.} None: the apparatus is parked for the
record; no graduating empirical claim was produced.

\paragraph{Key caveats from independent review.} The slow-lane
apparatus (three spec revisions, six verifier dispatches, a nine-cell
verdict cube) was overkill for the project's stage goal of identifying
a single empirically validated risk and a positive-\(\beta\) confirmation
before protocol design. The user pivoted to the simple-\(\beta\)
fast-lane track (Pair~D, dev-AI Stage-1) and parked P1 to preserve the
intellectual capital of the spec work. Re-activation conditions are
documented in the parking memory.

\paragraph{Pre-registration sha256 chain.} Spec sha256
\texttt{f855e036d3...b47aeab}.

\subsection{Cross-iteration multiple-comparison caveat}\label{sec:6.6}

The five iterations above are pre-registered individually but the
\emph{decisions} that produced them --- which target populations to
pursue, which \((Y, X)\) pairs to spec-pin, which to drop after an
EXIT, which to deepen --- are themselves a meta-multiple-comparison
structure. As of 2026-05-06 all five iterations are closed or parked
(Pair~D PASS; dev-AI Stage-1 FAIL; FX-vol-CPI FAIL; Phase-A.0
EXIT\_NON\_REMITTANCE; P1 SN18 PARKED). Pair~D's one-sided
\(p \approx 1.46 \times 10^{-8}\) is overwhelmingly significant against
any reasonable \citet{benjaminihochberg1995} or \citet{romanowolf2005}
correction applied across five iteration-level tests, but the
FX-vol-CPI pre-registered sensitivities (A1 monthly with \(90\%\)
interval \([+0.0057, +0.0246]\); A4 release-day-excluded with \(90\%\)
interval \([+0.0005, +0.0062]\)) and the dev-AI Stage-1 R2 Section-M
finding (\(\hat{\beta} = +0.455\), \(p = 1.13 \times 10^{-6}\), but
not pre-registered as a primary) are differentially robust to such a
correction; the A4 lower bound at \(0.0005\) is small in absolute
magnitude though the rejection is clean uncorrected. The Section-M
finding sits as a candidate-next-iteration object rather than a
graduated PASS, precisely because it was not pre-pinned as a primary
at spec-author time. We therefore decline to make any aggregate claim
about ``what the project has established across five iterations.''
What the project has established at the per-iteration level is
recorded above, iteration by iteration, with the appropriate hedges.
The framework's posture toward iteration-selection bias and the
appropriate familywise-error-rate or false-discovery-rate correction
regime is the dedicated open question we put to the supervisor in
Section~\ref{sec:7} (\(\S 7.5\)).

